\chapter*{Preface}
\addcontentsline{toc}{chapter}{Preface}

\section*{A book to read before \emph{and} after Strang}

There is a moment, somewhere in a first linear algebra course, when a student
writes down the equation $\mat{A}\vec{x} = \vec{b}$, solves it correctly, gets
full marks\ldots and still has no idea what just happened. The symbols moved
around the page in the legal ways. The answer matched the back of the book. But
the \emph{meaning} never arrived. Nothing in the calculation felt like it was
\emph{about} anything.

This book exists to make sure that does not happen to you. Or, if it has already
happened, to undo it.

Gilbert Strang's lectures and his \emph{Introduction to Linear Algebra} are,
deservedly, the way an enormous number of people have learned this subject. They
are warm, geometric, and full of the sense that linear algebra is one connected
idea rather than a list of tricks. If you have not met them yet, you have a
genuine pleasure ahead of you. So why another book?

Because Strang moves \emph{fast}, and he assumes you are ready to move with him.
He will draw a column picture and a row picture of a system of equations in the
first lecture, and if you have never before thought of a list of numbers as an
\emph{arrow}, or of a matrix as a \emph{machine that moves space around}, that
first lecture can go by like scenery from a train window: beautiful, gone before
you held it. This book is the slow walk through the same country \emph{before}
you take the train, and the deeper hike into the mountains \emph{after}.

\begin{intuitionbox}
The single picture this whole book is built to install in your head:
\begin{quote}
\textbf{A matrix is a function that takes in vectors and spits out vectors --
and it does so in the one geometrically simplest way possible: it keeps the grid
of space straight, evenly spaced, and pinned at the origin.}
\end{quote}
Everything else -- determinants, eigenvalues, rank, the SVD, least squares, the
math under modern machine learning -- is a consequence of that one sentence,
looked at from a different angle. If you finish a chapter and feel lost, come
back to this box. It is true on page 1 and it is true on the last page.
\end{intuitionbox}

\subsection*{Before Strang}

Read this book before Strang and its job is to build the intuition that makes his
course \emph{click} on the first pass instead of the third. When Strang says
``the column space is everything you can reach,'' you will already have spent an
afternoon reaching for things and noticing what you could and could not touch.
When he writes $\mat{A} = \mat{L}\mat{U}$, you will already feel elimination as a
sequence of physical shears on a grid, not as bookkeeping. You will walk into his
classroom not as a beginner but as someone who already has the mental furniture
and just needs the precise names for it.

\subsection*{After Strang}

Read this book after Strang and its job changes. Now you know the vocabulary, you
have solved the systems, you have computed the eigenvalues by hand. What you may
\emph{not} yet have is the feeling that the four fundamental subspaces, the
spectral theorem, and the singular value decomposition are the \emph{same theorem
wearing different clothes} -- and that this one theorem is precisely the thing
that powers principal component analysis, recommender systems, the stable
training of neural networks, and the compression of the image on your screen
right now. The later chapters of this book are written for the reader who is
ready to see the whole structure at once and to watch it turn into technology.

\begin{focusbox}
Hold onto two promises for the whole book:
\begin{enumerate}
\item \textbf{No formula appears before its picture.} You will see what a thing
\emph{means} before you are asked to compute it.
\item \textbf{No formula is asserted that we can derive in a few honest lines.}
When we can show you \emph{why}, we will. ``Trust me'' is not allowed here.
\end{enumerate}
\end{focusbox}

\section*{Who this book is for}

You are smart, you are motivated, and you may be seeing this material rigorously
for the first time. That is exactly the reader I am writing to. Concretely, this
book is for you if any of the following fit:

\begin{itemize}
\item You are about to take a linear algebra course (Strang's or any other) and
you want the geometry in your bones before the algebra starts.
\item You took such a course, passed it, and have the uncomfortable feeling that
you can compute things you cannot \emph{see}.
\item You came to linear algebra through machine learning, data science, computer
graphics, or signal processing, and you are tired of treating ``the matrix stuff''
as a black box you import from a library.
\item You simply think mathematics is one of the most beautiful things humans have
made, and you want to feel why this corner of it is so often called the most
useful mathematics ever invented.
\end{itemize}

\subsection*{What I assume you already have}

Very little, and all of it is the kind of thing you can refresh as you go.

\begin{itemize}
\item \textbf{High-school algebra.} You are comfortable with symbols standing for
numbers, with solving $2x + 1 = 7$, and with the idea of a function as a machine:
feed it an input, get an output.
\item \textbf{A little coordinate geometry.} You have seen the $xy$-plane, plotted
a point, drawn a line $y = mx + b$. We will rebuild everything else.
\item \textbf{Willingness to draw.} Not skill -- willingness. A shaky arrow on the
back of an envelope is worth more here than a paragraph of prose.
\end{itemize}

You do \emph{not} need calculus. Where calculus deepens a point, it lives in a
clearly marked box you can skip without losing the thread. You do not need any
prior programming, though if you can run a few lines of Python or use a
spreadsheet, some of the labs become more fun.

\begin{trapbox}
The most common way to misread this book is to treat the worked numerical
examples as something to \emph{watch} rather than something to \emph{do}. Reading
``$3 \cdot 2 + (-1)\cdot 4 = 2$'' and nodding feels like understanding. It is not.
Keep a pen moving. Recompute the small examples yourself, and when a box says
\emph{Compute it yourself}, actually stop and do it. Linear algebra is learned in
the hand before it is learned in the head.
\end{trapbox}

\section*{The philosophy: geometry, then algebra, then proof}

Almost every important idea in this book is introduced three times, in this order,
and the order is not decoration -- it is the whole pedagogical bet of the book.

\subsection*{First, geometry}

We start with a \emph{picture}. An arrow, a grid, a plane through the origin, a
square getting squashed into a parallelogram. The goal of the first pass is not
correctness in any formal sense; it is to give you something to \emph{look at} in
your mind's eye, so that the symbols which come later have somewhere to land. When
you can close your eyes and \emph{see} a matrix tilting the plane, you are no
longer memorizing -- you are describing something you can watch.

\subsection*{Then, algebra}

Once the picture is steady, we make it precise with coordinates and arithmetic. We
ask: if the picture is true, what must the \emph{numbers} do? This is where the
small, fully worked $2\times 2$ and $3\times 3$ examples live, traced line by
line, because the picture only becomes trustworthy once you have watched it agree
with the arithmetic in concrete cases.

\subsection*{Finally, proof}

A picture can mislead. A handful of examples can be a coincidence. So the third
pass asks the adult question: \emph{is it always true, and why?} Our proofs are
gentle but real. They are usually just the geometric idea written carefully enough
that a skeptic could not wriggle out. You will find that, done this way, a proof
is not a hurdle after the understanding -- it \emph{is} the understanding,
pinned down so it cannot drift.

\begin{intuitionbox}
A tiny illustration of the three passes, so the philosophy is not just a slogan.
Take the claim: \emph{rotating the plane by $90^\circ$ never changes a vector's
length.}

\textbf{Geometry.} Spin an arrow about the origin like the hand of a clock. It
points somewhere new; it is plainly just as long. Length cannot care which way
you face.

\textbf{Algebra.} The rotation sends $\colvec{x \\ y}$ to $\colvec{-y \\ x}$.
Its squared length is $(-y)^2 + x^2 = x^2 + y^2$ -- the same number we started
with. Try it on $\colvec{3 \\ 1}$: it goes to $\colvec{-1 \\ 3}$, and
$9 + 1 = 1 + 9 = 10$. Length preserved.

\textbf{Proof.} For \emph{every} $\colvec{x \\ y}$ the swap-and-negate sends
$x^2 + y^2$ to $(-y)^2 + x^2$, and addition does not care about order, so the two
sums are identical. Always. No exceptions. The picture was telling the truth.
\end{intuitionbox}

\begin{walkbox}
Before you read another page, take this on a walk. You have two arrows lying flat
on a table, starting from the same point. You are allowed to \emph{add} them
(tip to tail) and to \emph{scale} them (stretch or flip). Which points on the
table can you reach by adding and scaling? When can you reach \emph{every} point,
and when are you stuck on a single line forever? You will have just discovered,
on your own, the central drama of the first third of this book: \emph{span},
\emph{independence}, and \emph{dimension}. Notice how far you get with no
formula at all.
\end{walkbox}

\section*{The colored boxes: a visual language}

This book talks to you in two registers. The main text walks forward steadily.
Around it sit \emph{colored boxes}, each a different kind of voice -- a nudge, a
warning, a moment of beauty, a bridge to the wider world. They are not optional
sidebars to skim. They are where a lot of the real learning happens, because they
are where the book stops marching and turns to look you in the eye. Here is the
full cast, so you know who is speaking the moment you see a color.

\begin{intuitionbox}
\textbf{Intuition (green).} The core mental picture. When you see green, slow
down and try to \emph{see} the thing being described before you read on. If you
remember nothing else from a section, remember its green boxes.
\end{intuitionbox}

\begin{focusbox}
\textbf{Focus (blue).} The one or two sentences worth holding onto from a stretch
of text. If you were making flashcards, these are the fronts of the cards.
\end{focusbox}

\begin{trapbox}
\textbf{Common trap (red).} A classic mistake -- one that almost everyone makes,
including people who go on to use this subject professionally. Seeing the trap
named in advance is often what keeps you out of it.
\end{trapbox}

\begin{walkbox}
\textbf{Walking question (orange).} A prompt to chew on away from the page,
ideally while doing something else. The best linear algebra thinking I have ever
done happened while walking, not while staring. These questions have no
deadline; let them simmer.
\end{walkbox}

\begin{labbox}
\textbf{Compute it yourself (purple).} A small, do-it-by-hand $2\times2$ or
$3\times3$ lab. Pen, paper, two minutes. These are the reps that turn watching
into knowing.
\end{labbox}

\begin{formalbox}
\textbf{Formal (gray).} The precise statement -- a definition or theorem stated
exactly, the way you would need it to prove something or to look it up later.
When you want to be sure of the \emph{letter} of an idea, find its gray box.
\end{formalbox}

\begin{strangbox}
\textbf{Before/after Strang (deep blue).} An explicit bridge to how Strang frames
the same idea. \emph{Before} you have taken his course, these boxes tell you what
to listen for. \emph{After}, they connect his vocabulary -- pivots, column space,
the four subspaces -- to the way we have built things here, so the two pictures
fuse into one. Example: where we will speak of ``the directions a matrix flattens
to zero,'' Strang calls that the \emph{null space}, and you will hear him say
``$\mat{A}\vec{x}=\vec{0}$'' for exactly the same idea.
\end{strangbox}

\begin{applicationbox}
\textbf{Modern application (magenta).} Where this idea shows up in the world right
now -- in machine learning, data science, computer graphics, signals, networks,
or statistics. A taste of why people sometimes call linear algebra ``the language
of data.'' For instance: every time a streaming service guesses a film you will
like, it is approximating a giant matrix of ratings by a product of two thin ones
-- a low-rank factorization, which is a singular value decomposition wearing a
business suit. We will earn that sentence completely by the end of the book.
\end{applicationbox}

\begin{connectionbox}
\textbf{Connection to other math (teal).} A bridge to a neighboring country:
calculus, geometry, group theory, probability, differential equations, functional
analysis. Linear algebra is the hub that these spokes meet at, and seeing the
spokes is part of seeing why it matters so much. For example: the derivative
you met in calculus is, at heart, the \emph{best linear approximation} to a curve
at a point -- which is to say, calculus is the art of replacing complicated
functions by matrices, locally, and linear algebra is what you do with those
matrices once you have them.
\end{connectionbox}

\begin{aestheticbox}
\textbf{Why this matters (cyan).} A pause to say plainly why a result is beautiful
or important -- the part a hurried course skips, and the part that, for many
people, is the whole reason to bother. Mathematics is not only useful; it is one
of the most reliably gorgeous things people make, and I will not pretend
otherwise.
\end{aestheticbox}

\begin{historybox}
\textbf{History and connection (yellow).} A little context on where an idea came
from and who fought to understand it -- because the ideas were hard-won by real
people, and knowing that they were once confused too is strangely encouraging.
\end{historybox}

\begin{recoverybox}
\textbf{Pause (plain).} A breath. A short recap, or permission to stop, reread,
and let things settle before going on. If a chapter ever feels like too much at
once, the Pause boxes are designed to be safe places to rest.
\end{recoverybox}

\begin{marginnotebox}
\textbf{Reading note (gray).} A small piece of advice about \emph{how} to read a
particular passage -- what to skim, what to linger on, what is safe to skip on a
first pass and return to later.
\end{marginnotebox}

\begin{focusbox}
You do not need to memorize this list. The colors are chosen so that, within a
chapter or two, you will know what kind of voice is speaking before you have read
a single word -- green means \emph{see this}, red means \emph{careful}, magenta
means \emph{the real world}, teal means \emph{the rest of math}. The page itself
will start telling you how to read it.
\end{focusbox}

\section*{Pictures are first-class citizens here}

In many textbooks the diagrams are illustrations: pleasant, optional, decorative.
In this one they are \emph{arguments}. A figure that shows the unit square getting
sheared into a parallelogram is not prettifying the determinant -- it \emph{is}
the determinant, which is nothing more than the factor by which area got scaled.
You will see, drawn carefully and referenced from the prose:

\begin{itemize}
\item vectors as arrows, and their sums as journeys laid tip to tail;
\item linear combinations and \emph{spans} -- the lines and planes you can reach;
\item the action of a matrix on a whole grid, and on the humble unit square;
\item \emph{projections} dropping a point straight down onto a line or a plane;
\item subspaces nested inside one another, and the famous four of them;
\item eigenvector directions -- the special arrows a matrix does not turn;
\item rotations, reflections, and shears, each with its own visual fingerprint;
\item and, as the book's quiet destination, the singular value decomposition as
\emph{rotate, stretch, rotate} -- the claim that \emph{every} matrix, no matter
how complicated, is only ever those three simple moves in sequence.
\end{itemize}

Below is the kind of small, honest figure you will meet on nearly every page.
Nothing fancy: just an arrow, its double, and their sum, so that the very first
operations of linear algebra -- scaling and adding -- are things you have
\emph{seen} and not only read.

\begin{figure}[H]
\centering
\begin{tikzpicture}[scale=1.0,>=Stealth]
  % light grid
  \draw[step=1,gray!22,very thin] (-0.4,-0.4) grid (5.4,3.4);
  % axes
  \draw[->,gray!60] (-0.4,0) -- (5.4,0) node[right]{$x$};
  \draw[->,gray!60] (0,-0.4) -- (0,3.4) node[above]{$y$};
  % vector v
  \draw[->,very thick,vecblue] (0,0) -- (2,1)
        node[midway,above left,vecblue]{$\vec{v}$};
  % 2v
  \draw[->,very thick,vecblue!55] (0,0) -- (4,2)
        node[midway,below right,vecblue!70]{$2\vec{v}$};
  % vector w drawn from origin (faint) ...
  \draw[->,very thick,vecred] (0,0) -- (1,2)
        node[midway,above left,vecred]{$\vec{w}$};
  % ... and w again from the tip of v, to show the sum journey
  \draw[->,very thick,vecred!55,dashed] (2,1) -- (3,3);
  % the sum v + w
  \draw[->,ultra thick,structurepurple] (0,0) -- (3,3)
        node[midway,above,structurepurple]{$\vec{v}+\vec{w}$};
  \fill (0,0) circle (1.6pt);
\end{tikzpicture}
\caption{Two operations are the whole foundation. Scaling: $\vec{v}=\colvec{2\\1}$
stretches to $2\vec{v}=\colvec{4\\2}$ along the same line. Adding: walk along
$\vec{v}$, then continue along $\vec{w}=\colvec{1\\2}$ (dashed), and you arrive at
$\vec{v}+\vec{w}=\colvec{3\\3}$. Every span, every matrix, every theorem in this
book is built out of just these two moves.}
\label{ch00:scale-and-add}
\end{figure}

Figure~\ref{ch00:scale-and-add} is doing real work, not decoration. The dashed
copy of $\vec{w}$ \emph{is} the reason vector addition is ``tip to tail''; the two
shades of blue \emph{are} the reason scaling stays on one line. Whenever a figure
appears, slow down and ask what claim it is quietly making.

Here is a second taste of a figure-as-argument -- one we will return to so often
that it is almost the mascot of the book. Take the humble unit square, the one
cornered at the origin with sides $\colvec{1\\0}$ and $\colvec{0\\1}$, and feed it
to the matrix $\mat{A}=\colvec{1 & 1 \\ 0 & 1}$. The matrix sends each corner to a
new place, and the square slides into a \emph{parallelogram}. Watch what stays
true and what changes.

\begin{figure}[H]
\centering
\begin{tikzpicture}[scale=1.1,>=Stealth]
  % left panel: unit square
  \begin{scope}[shift={(0,0)}]
    \draw[step=1,gray!20,very thin] (-0.3,-0.3) grid (2.3,2.3);
    \draw[->,gray!60] (-0.3,0) -- (2.3,0) node[right]{$x$};
    \draw[->,gray!60] (0,-0.3) -- (0,2.3) node[above]{$y$};
    \fill[vecblue!18] (0,0) -- (1,0) -- (1,1) -- (0,1) -- cycle;
    \draw[very thick,vecblue] (0,0) -- (1,0) -- (1,1) -- (0,1) -- cycle;
    \draw[->,very thick,vecred] (0,0) -- (1,0);
    \draw[->,very thick,structurepurple] (0,0) -- (0,1);
    \node at (0.5,-0.7) {\small area $=1$};
  \end{scope}
  % arrow between
  \draw[->,very thick,gray!70] (2.7,1) -- (3.9,1)
        node[midway,above]{$\mat{A}$};
  % right panel: sheared square
  \begin{scope}[shift={(4.3,0)}]
    \draw[step=1,gray!20,very thin] (-0.3,-0.3) grid (3.3,2.3);
    \draw[->,gray!60] (-0.3,0) -- (3.3,0) node[right]{$x$};
    \draw[->,gray!60] (0,-0.3) -- (0,2.3) node[above]{$y$};
    \fill[vecblue!18] (0,0) -- (1,0) -- (2,1) -- (1,1) -- cycle;
    \draw[very thick,vecblue] (0,0) -- (1,0) -- (2,1) -- (1,1) -- cycle;
    \draw[->,very thick,vecred] (0,0) -- (1,0);
    \draw[->,very thick,structurepurple] (0,0) -- (1,1);
    \node at (1.0,-0.7) {\small area $=1$ still};
  \end{scope}
\end{tikzpicture}
\caption{A matrix is a motion of space. The shear $\mat{A}=\protect\colvec{1 & 1
\\ 0 & 1}$ leaves the red side $\protect\colvec{1\\0}$ alone and tilts the purple
side $\protect\colvec{0\\1}$ over to $\protect\colvec{1\\1}$. The square becomes a
slanted parallelogram, yet its area is unchanged: base $1$, height $1$, area $1$.
That unchanged area is the \emph{determinant}, here equal to $1$. We have not
defined the determinant yet -- but you have already seen it.}
\label{ch00:shear-square}
\end{figure}

Figure~\ref{ch00:shear-square} is the entire idea of a determinant, smuggled in
before its name. The lesson of the whole book is hiding in that pair of pictures:
\emph{compute nothing, see everything first.} When the formula
$\det\protect\colvec{a & b \\ c & d} = ad - bc$ finally arrives, it will be the
caption to a figure you already trust, not a rule to memorize.

\begin{labbox}
A thirty-second warm-up, to prove to yourself that the figure is honest. On the
same axes, draw $\vec{v} = \colvec{2 \\ 1}$ and $\vec{w} = \colvec{1 \\ 2}$, then
form $\vec{v}+\vec{w}$ two ways: once by walking $\vec{v}$ then $\vec{w}$, and once
by walking $\vec{w}$ then $\vec{v}$. You should land on $\colvec{3\\3}$ both times.
You have just \emph{seen} that vector addition does not care about order -- a fact
we will lean on constantly, and which you now believe because you watched it,
not because anyone told you.
\end{labbox}

\section*{How the parts build}

The book climbs in deliberate stages. You can read it straight through, and that
is what I recommend on a first pass, because each part quietly installs the tools
the next one stands on. Here is the arc, so you can see the mountain before you
start up it.

\begin{enumerate}
\item \textbf{Vectors and the spaces they live in.} We begin where everything
begins: a vector as an arrow and as a list of numbers, the two operations of
scaling and adding, and the geometry of \emph{span} -- the lines, planes, and
spaces you can build. This is where ``before Strang'' readers get the furniture
that the column and row pictures will later sit on.

\item \textbf{Matrices as motions.} A matrix stops being a grid of numbers and
becomes a \emph{transformation} -- a way of moving all of space at once while
keeping the grid straight. Matrix multiplication becomes \emph{composition} of
motions, and suddenly the strange ``row times column'' rule is the only rule it
could possibly have been.

\item \textbf{Solving and structure.} Systems of equations, elimination, rank,
and the four fundamental subspaces -- the skeleton Strang builds his whole course
around. We earn each piece geometrically first, so the algebra of elimination
feels like watching a shape simplify rather than pushing symbols.

\item \textbf{Geometry made precise.} Length, angle, the dot product,
\emph{orthogonality}, and \emph{projection} -- the tools that let us measure, and
then the crown jewel of applied linear algebra, \emph{least squares}: the best
answer when there is no exact answer, which is the situation nearly all real data
puts you in.

\item \textbf{Eigenvalues and the heart of a matrix.} The special directions a
matrix merely stretches without turning, why they reveal a matrix's deepest
behavior, and how they let us understand repeated processes -- powers, dynamics,
the long-run fate of a system.

\item \textbf{The summit: the singular value decomposition.} Every matrix as
\emph{rotate, stretch, rotate}. From here we can finally see, fully and honestly,
the machinery under principal component analysis, data compression, and a great
deal of modern machine learning. This is the chapter the whole book was walking
toward.
\end{enumerate}

\begin{strangbox}
If you are reading this \emph{before} Strang: parts 1 through 3 are your runway.
By the time you have them, his opening lectures will feel like meeting old friends
under new names. If you are reading \emph{after} Strang: feel free to skim the
early parts for the geometric pictures you may have rushed past, and spend your
real energy in parts 4 through 6, where we go deeper than a first course usually
can -- especially into the SVD, which Strang rightly calls the high point, and
which here gets the slow, full treatment it deserves.
\end{strangbox}

\begin{aestheticbox}
Here is the promise that the structure above is really making. By the final part,
six ideas that looked separate at the start -- linear combinations, matrix
multiplication, rank, orthogonal projection, eigenvalues, and singular values --
will have collapsed into \emph{one} idea seen from six sides. That collapse, the
moment a sprawling subject suddenly closes into a single clean shape in your mind,
is one of the great pleasures mathematics has to offer. It is worth the climb. I
have tried to build every step so that you actually get to feel it.
\end{aestheticbox}

Two of those six ideas deserve a face before you meet them by name, because they
are the ones students most often compute without ever seeing. Figure~\ref{ch00:proj-eigen}
shows them side by side: a \emph{projection} on the left, an \emph{eigenvector} on
the right.

\begin{figure}[H]
\centering
\begin{tikzpicture}[scale=1.0,>=Stealth]
  % --- left: projection onto a line ---
  \begin{scope}[shift={(0,0)}]
    \draw[->,gray!55] (-0.3,0) -- (3.6,0) node[right]{$x$};
    \draw[->,gray!55] (0,-0.3) -- (0,3.2) node[above]{$y$};
    % the line we project onto: direction (3,1)
    \draw[thick,spacegray] (-0.2,-0.0667) -- (3.6,1.2)
          node[right,spacegray]{\small $\ell$};
    % the point / vector b
    \coordinate (b) at (1,2.6);
    \draw[->,very thick,vecred] (0,0) -- (b) node[above left,vecred]{$\vec{b}$};
    % projection of b onto line: p = ((b.d)/(d.d)) d, d=(3,1)
    % b.d = 3+2.6=5.6 ; d.d=10 ; t=0.56 ; p=(1.68,0.56)
    \coordinate (p) at (1.68,0.56);
    \draw[->,very thick,vecblue] (0,0) -- (p)
          node[below right,vecblue]{\small $\hat{\vec{b}}$};
    % the dropped perpendicular (the error)
    \draw[dashed,thick,structurepurple] (b) -- (p);
    % right-angle mark
    \draw (p) ++(-0.18,0.06) -- ++(0.06,0.18) -- ++(0.18,-0.06);
    \node at (1.9,2.3) {\small \color{structurepurple}error};
    \node[align=center] at (1.7,-0.7) {\small projection onto a line};
  \end{scope}
  % --- right: eigenvector ---
  \begin{scope}[shift={(5.6,0)}]
    \draw[->,gray!55] (-0.3,0) -- (3.6,0) node[right]{$x$};
    \draw[->,gray!55] (0,-0.3) -- (0,3.2) node[above]{$y$};
    % A=[[1.5,0.5],[0.5,1.5]] has eigvec (1,1) with eigval 2.
    % v=(1.3,1.3) -> Av = 2v = (2.6,2.6).
    \draw[->,very thick,vecblue] (0,0) -- (1.3,1.3)
          node[above left,vecblue]{\small $\vec{v}$};
    \draw[->,ultra thick,structurepurple] (0,0) -- (2.6,2.6)
          node[right,structurepurple]{\small $\mat{A}\vec{v}=2\vec{v}$};
    % a non-eigen vector gets turned. u=(2,0.4) -> Au=(3.2,1.6) exactly.
    \draw[->,very thick,vecred!70] (0,0) -- (2.0,0.4)
          node[below right,vecred]{\small $\vec{u}$};
    \draw[->,very thick,vecred!40,dashed] (0,0) -- (3.2,1.6)
          node[above right,vecred!70]{\small $\mat{A}\vec{u}$};
    \node[align=center] at (1.8,-0.7) {\small the arrow $\mat{A}$ won't turn};
  \end{scope}
\end{tikzpicture}
\caption{Two pictures the book earns slowly. \textbf{Left:} the projection
$\hat{\vec{b}}$ is the shadow of $\vec{b}$ cast straight down onto the line
$\ell$; the purple dashed segment is the \emph{error}, and it meets $\ell$ at a
right angle -- that perpendicularity is the entire secret of least squares.
\textbf{Right:} most arrows get turned by a matrix (red $\vec{u}\to\mat{A}\vec{u}$),
but a special one, the \emph{eigenvector} $\vec{v}$, is only stretched -- here
$\mat{A}\vec{v}=2\vec{v}$, doubled but not rotated.}
\label{ch00:proj-eigen}
\end{figure}

\begin{labbox}
A two-minute numeric preview, so the words above have a number behind them. Let
$\mat{A}=\colvec{1.5 & 0.5 \\ 0.5 & 1.5}$ and $\vec{v}=\colvec{1\\1}$. Multiply,
row by row: the top entry is $1.5\cdot 1 + 0.5\cdot 1 = 2$, and the bottom entry
is $0.5\cdot 1 + 1.5\cdot 1 = 2$. So $\mat{A}\vec{v}=\colvec{2\\2}=2\vec{v}$: the
arrow doubled and did not turn. Now try $\vec{u}=\colvec{2\\0}$ and watch it get
\emph{bent} instead: $\mat{A}\vec{u}=\colvec{3\\1}$, which no longer points along
the $x$-axis. You have just told an eigenvector from an ordinary vector by hand --
the entire fifth part of the book is the slow, careful version of this one
calculation.
\end{labbox}

\section*{How to read this book}

A few practical suggestions, earned from watching people learn this material.

\begin{itemize}
\item \textbf{Keep a pen in your hand.} I will say it once more because it is the
single highest-leverage habit: redo the small examples, attempt the purple labs,
and sketch the figures yourself. Passive reading produces the illusion of
understanding and very little of the real thing.

\item \textbf{Let the orange questions breathe.} The walking questions are not
graded and not urgent. Carry one around for a day. The understanding it produces
is deeper than anything you can get by reading the next paragraph faster.

\item \textbf{When you are lost, go back to green.} If a section dissolves into
fog, find the nearest Intuition box and the nearest figure, and rebuild the
picture before pushing on. Almost all confusion in this subject is a picture that
slipped, not a calculation that broke.

\item \textbf{Do the exercises.} Every chapter ends with a graded set, from
warm-ups to genuine challenges, mixing computation, geometry, and a proof or two.
They are where the chapter's ideas stop being things you have read and become
things you can use. You do not need to do all of them, but do not skip them
wholesale -- pick a warm-up, a geometry problem, and a hard one from each
chapter at minimum.

\item \textbf{It is fine to be confused for a while.} Confusion is not failure;
it is the feeling of your mind reshaping itself around a genuinely new idea.
Every mathematician you admire has sat exactly where you are sitting, stuck on
exactly the thing you are stuck on. Stay with it. Take a walk. Come back. The
picture will steady.
\end{itemize}

\begin{historybox}
Take heart from the history. The objects in this book did not arrive clean. People
solved systems of equations for two thousand years before anyone wrote a
\emph{matrix}; the word itself is younger than the steam engine. Determinants were
studied for a century before the matrices they belong to. Eigenvalues were buried
inside the study of vibrating strings and spinning planets long before they had a
name. The tidy, unified subject you are about to learn is the polished result of
centuries of confusion exactly like the confusion you will sometimes feel. You are
in excellent company.
\end{historybox}

\begin{connectionbox}
One last reason to keep going, for the reader who wonders where this leads. The
``keep the grid straight'' idea at the center of this book does not stop at arrows
in the plane. Replace arrows with \emph{functions} and the same machinery becomes
\emph{functional analysis}, the language of quantum mechanics and of the Fourier
transform that lives inside every audio and image codec. Replace addition with
\emph{composition} and you brush up against \emph{group theory}, the mathematics
of symmetry. Linear algebra is not a corner of mathematics; it is a crossroads,
and almost every road worth walking passes through it. This book is your map of
the crossroads.
\end{connectionbox}

\section*{Warm-up exercises}

A preface has no theorems to drill, but it can still ask you to pick up a pen --
which, as you now know, is the whole point. None of these needs anything past
high-school algebra and the two figures above. Treat them as a gentle handshake
with the book's habits: \emph{draw first, compute second, and trust the picture
once the numbers agree with it.}

\begin{enumerate}
\item \textbf{(Tip to tail.)} On graph paper draw $\vec{v}=\colvec{2\\1}$ and
$\vec{w}=\colvec{-1\\2}$ from the origin. Form $\vec{v}+\vec{w}$ by the tip-to-tail
rule, read its coordinates off the page, and then check by adding componentwise.
Do the two answers match? They had better.

\item \textbf{(Scaling stays on a line.)} With the same $\vec{v}=\colvec{2\\1}$,
plot $\tfrac{1}{2}\vec{v}$, $2\vec{v}$, and $-\vec{v}$. Convince yourself by eye
that all four tips lie on a single line through the origin. Which one points the
opposite way, and why?

\item \textbf{(Rotation preserves length, again.)} The text claimed a $90^\circ$
turn sends $\colvec{x\\y}$ to $\colvec{-y\\x}$. Apply it to $\colvec{4\\3}$, then
to the result, then twice more. After four turns, where are you? Compute the
length $\sqrt{x^2+y^2}$ at each step and confirm it never changes.

\item \textbf{(A shear, by hand.)} Using $\mat{A}=\colvec{1 & 1 \\ 0 & 1}$ from
Figure~\ref{ch00:shear-square}, compute where the four corners of the unit square
$\colvec{0\\0},\colvec{1\\0},\colvec{1\\1},\colvec{0\\1}$ land. Plot the image.
You should recover the slanted parallelogram of the figure -- and you can check
its area is still $1$ by counting base times height.

\item \textbf{(Spot the eigenvector.)} For $\mat{A}=\colvec{1.5 & 0.5 \\ 0.5 &
1.5}$, the purple lab found that $\colvec{1\\1}$ is merely stretched. Test the
candidate $\colvec{1\\-1}$ the same way: compute $\mat{A}\colvec{1\\-1}$. Is it a
scalar multiple of $\colvec{1\\-1}$? If so, by what factor? (You will have found
this matrix's \emph{second} eigenvector without any of the machinery -- a small
preview of how clean Part~5 turns out to be.)

\item \textbf{(A walking question, on paper for once.)} You have two arrows in
the plane, not pointing along the same line. Argue in a sentence or two why you
can reach \emph{every} point of the plane by scaling and adding them -- and what
goes wrong the moment the two arrows \emph{do} share a line. This is the seed of
\emph{span} and \emph{independence}; say it in your own words now, and watch the
book sharpen your wording later.
\end{enumerate}

\section*{A note of welcome}

Linear algebra has a reputation for being dry -- a gauntlet of row reductions to
survive on the way to something more interesting. I think that reputation is a
tragedy of bad first encounters, and this book is my attempt to give you a good
one. Underneath the row reductions is a subject of startling beauty and almost
unreasonable usefulness: the mathematics that lets us understand data, render
worlds, model the physical universe, and -- right now, in our moment -- build
machines that learn. You are about to see how a single, simple idea about
keeping space straight blossoms into all of that.

So sharpen a pencil. Find some scrap paper. We are going to start, quite
literally, from the ground up -- with a single arrow, drawn from the origin, and
the two simplest things you can do to it.

Welcome. Let us begin.

\begin{flushright}
\emph{Turn the page, and meet a vector.}
\end{flushright}
